% Pursuing Stacks (cotes 134-2 to 134-8): archive cote 134-7, site page 112 (PDF page 113), typescript page 554.
% typescript. Draft compared: ps6.tex:5645-5703.
% AI-transcribed by Claude Opus 5.5 (claude-opus-5-5, adaptive thinking, effort high), Message Batches API, batch ps4 (msgbatch_012Aoc1NuskATPsyBeKSYK7J), 2026-11-12; not yet checked by eye.
\scanpage{134-7/112}
\mspage{554}

the would-be expression of ``total co-homotopy'' of $X_*$ we'll get this
way is'nt so\xout{x} much more appealing than (18) -- it lies still in one of
the wrong quadrants, which definitely makes us feel unconfortable.

In principle, the tentative formula (18),\xout{gives} when applied say
to an object such as
\[ S(n,k)\overset{\mathrm{def}}{=} \overset{\hand{\unsure{$\sim$}}}{S}(\, k[n]) \quad , \]
gives a rather explicit (but for the time being highly hypothetical\hand{!})
expression of the homotopy modules of the n-sp\over{\unsure{ae}}{he}re over k, which in
case $k=\underline{Z}$ are hoped to be just the homotopy groups of the ordinary
n-sphere. To test whether this makes at all a sense, we'll have to
understand first the structure of the \xout{Lie} André-Quillen ``Lie complex''
of an algebra \inserted{(6)} of the type D(M)\hand{,} \over{\unsure{f}}{f}or variable k-module M. I hav'nt started
looking into this yet, and I doubt I am going to do it presently.

At any rate, whether or not the formula (18) we ended up with
is essentially correct, in order \inserted{(among other things)} to get a \inserted[typed]{method} \xout{principle} for computing
the total homotopy $L\hand{\pi}_{\bullet}$ for a semisimplicial bundle $X_*$ which is'nt a
flat vector bundle, we'll have to find out some more or less explicit
means of replacing $X_*$ by some $X'_*$ \xout{which} whose components \uline{are} flat vector
bundles, and which is isomorphic to $X_*$ in the relevant derived category
$\mathrm{Hot}_o(k)$. When k is'nt a fi\over{\unsure{ed}}{el}d, the question arises already even when \margin{field}
\xout{$X_*$} the components $X_n$ \uline{are} vector bundles, when these are not flat. T\over{\unsure{h}}{h}e
first idea that comes to mind, from the André-Quillen theory precisely,
\xout{is that to} in terms of the co-semisimplicial algebra $A^*$ expressing $X_*$ ,
is to take ``projective resolutions'' of \xout{each xxxx} the various components
$A^n$ by polynomial algebras. Again we end up with a mixed complex, this
time a complex of augmented algebras depending on two indices n,p ,
covariant in n and contravariant in p, or \xout{xxxxx} the reverse if we
replace those algebras $A^n_p$ by their \unsure{\uline{spectra}} $X_n^p$ . In any case,
we again end up in a ``wrong quadrant''! \margin{$\|$}

The hesitating question that
comes to mind now is whether it is at all feasible to work with a catego-%
ry of models which is'nt a category of \uline{semisimplicial} bundles say, but
\inserted{one} of such \uline{mixed} \uline{wrong-quadrant\deleted{s}} bundles\over{,}{;} \deleted{and} \inserted{namely,} use these as ``models'' for
getting hold of a reasonable derived category of ``schematic homotopy
types''\hand{?} I never heard of anything such yet, and I confess that at this
point my \inserted{(anyhow rather poor!)} formal intuition of the situation breaks down completely --
maybe the suggestion is complete nonsense, for some wholly trivial
reason ! Maybe Larry Breen could tell me at once -- or someone else
who has more feeling than I for semisimplicial and cosemisimplicial mode\unsure{ls}

\margin{112}
